% Standalone problem document, generated from the adjacent README.md.
% Compile with XeLaTeX. Mathematical content is maintained in Markdown.
\documentclass[11pt,a4paper]{article}
\usepackage[fontsize=10.5pt]{scrextend}
\usepackage[margin=22mm,headheight=15pt,headsep=9mm,footskip=11mm]{geometry}
\usepackage{fontspec}
\setmainfont{texgyrepagella}[Extension=.otf,UprightFont=*-regular,BoldFont=*-bold,ItalicFont=*-italic,BoldItalicFont=*-bolditalic]
\setsansfont{texgyreheros}[Extension=.otf,UprightFont=*-regular,BoldFont=*-bold,ItalicFont=*-italic,BoldItalicFont=*-bolditalic]
\setmonofont{texgyrecursor}[Extension=.otf,UprightFont=*-regular,BoldFont=*-bold,ItalicFont=*-italic,BoldItalicFont=*-bolditalic,Scale=MatchLowercase]
\usepackage{amsmath,amssymb}
\usepackage{unicode-math}
\setmathfont{texgyrepagella-math.otf}
\usepackage{xcolor,microtype,enumitem,needspace,fancyhdr}
\usepackage{bookmark}
\definecolor{ink}{HTML}{17344B}
\definecolor{accent}{HTML}{197D80}
\definecolor{muted}{HTML}{596773}
\hypersetup{colorlinks=true,linkcolor=accent,urlcolor=accent,pdftitle={SP-04 - The
smallest-multiplier rule for nearest unit-absolute-determinant
matrices},pdfauthor={Open Problems in Numerical Linear Algebra}}
\urlstyle{same}
\setlength{\parindent}{0pt}
\setlength{\parskip}{5pt plus 1pt minus 1pt}
\linespread{1.04}
\setlength{\emergencystretch}{3em}
\widowpenalty=10000
\clubpenalty=10000
\displaywidowpenalty=10000
\allowdisplaybreaks[1]
\setlist{nosep,leftmargin=1.5em}
\providecommand{\tightlist}{\setlength{\itemsep}{0pt}\setlength{\parskip}{0pt}}
\providecommand{\pandocbounded}[1]{#1}
\pagestyle{fancy}
\fancyhf{}
\fancyhead[L]{\sffamily\footnotesize\color{muted}OPEN PROBLEMS IN NUMERICAL LINEAR ALGEBRA}
\fancyhead[R]{\sffamily\bfseries\footnotesize\color{accent}SP-04}
\fancyfoot[L]{\parbox[b]{0.9\textwidth}{\sffamily\fontsize{8}{10}\selectfont\color{muted}Editorial ratings; a bounded search cannot certify that a problem remains unsolved.\\Literature check: 2026-09-10}}
\fancyfoot[R]{\sffamily\footnotesize\color{muted}\thepage}
\renewcommand{\headrulewidth}{0.3pt}
\renewcommand{\headrule}{\hbox to\headwidth{\color{accent}\leaders\hrule height \headrulewidth\hfill}}
\makeatletter
\renewcommand{\section}{\Needspace{5\baselineskip}\@startsection{section}{1}{0pt}{12pt plus 2pt minus 2pt}{4pt}{\sffamily\bfseries\large\color{ink}}}
\renewcommand{\subsection}{\Needspace{4\baselineskip}\@startsection{subsection}{2}{0pt}{9pt plus 1pt minus 1pt}{3pt}{\sffamily\bfseries\normalsize\color{ink}}}
\makeatother
\setcounter{secnumdepth}{0}
\begin{document}
{\sffamily\small\color{muted}Eigenvalues and inverse problems\par}
\vspace{3pt}
{\raggedright\hyphenpenalty=10000\exhyphenpenalty=10000\sffamily\bfseries\fontsize{21}{24}\selectfont\color{ink}The
smallest-multiplier rule for nearest unit-absolute-determinant
matrices\par}
\vspace{7pt}
{\sffamily\small\textbf{Difficulty:} challenging\quad\textbf{Importance:} interesting
to specialist\par}
{\sffamily\small\bfseries\color{accent}Status: Open\par}
\vspace{4pt}
{\color{accent}\hrule height 0.8pt}
\vspace{6pt}
\textbf{Rating rationale:} Challenging reflects a global root-selection
rule across all dimensions and singular-value data; specialist impact is
the exact projection onto a specific determinant constraint.

For \(n\ge2\), define \[
\mathcal G_n=\{X\in\mathbb R^{n\times n}:\det X\in\{-1,1\}\},
\qquad
\|Y\|_F^2=\sum_{i,j}y_{ij}^2.
\] For a real data matrix \(U\), consider the real stationary pairs \[
\mathcal C(U)=
\{(X,c)\in\mathcal G_n\times\mathbb R:X^T(U-X)=cI_n\}.
\] For generic \(U\), this set is finite. Choose a pair \((X_*,c_*)\)
whose multiplier has minimum absolute value. Does \[
\|U-X_*\|_F=\min_{X\in\mathcal G_n}\|U-X\|_F
\] hold for every \(n\ge2\) and generic real \(U\)?

Here generic means outside an exceptional proper real algebraic set; in
particular \(U\) is invertible and its squared singular values are
distinct. If minimum absolute multipliers tie, the question can be
restricted to the generic locus where the choice is unique. The source
phrases the same selection question using the real roots of its
elimination polynomial \(R_1(c)\), which correspond one-to-one to
stationary matrices for generic data.

\subsection{Relevance}\label{relevance}

This is a concrete question about selecting the global nearest matrix
from the stationary solutions of a structured matrix nearness problem. A
positive answer would supply a scalar criterion for identifying the
global solution. Both determinant signs are allowed; this must not be
silently replaced by projection onto \(\operatorname{SL}_n(\mathbb R)\)
alone.

\newpage
\subsection{References}\label{references}

\begin{itemize}
\tightlist
\item
  J. A. Baaijens and J. Draisma, \emph{Euclidean distance degrees of
  real algebraic groups}, Linear Algebra and its Applications 467
  (2015), 174--187, §4.1, equations (4)--(5), elimination algorithm, and
  Problem 4.3 on p.186
  (\href{https://pure.tue.nl/ws/files/3846347/391917266748824.pdf}{version
  of record};
  \href{https://doi.org/10.1016/j.laa.2014.11.012}{journal}).
\item
  P. Jaap and O. Sander, \emph{How to project onto SL(n)},
  arXiv:2501.19310 (2025), introduction pp.1--2, discussing the
  Baaijens--Draisma conjectural global selection rule
  (\href{https://arxiv.org/pdf/2501.19310}{primary manuscript}); Journal
  of Optimization Theory and Applications 209 (2026), article 40
  (\href{https://doi.org/10.1007/s10957-026-02974-8}{journal}).
\end{itemize}

\subsection{Status check --- 2026-09-08}\label{status-check-2026-09-08}

The January 2025 manuscript explicitly describes the earlier proposed
global-minimum selection as conjectural. Its journal version appeared in
April 2026; the accessible journal abstract and arXiv version history
were checked, but the subscription-only journal body was not inspected.
Searches for ``nearest determinant-one matrix'', ``Baaijens'',
``smallest absolute value'', ``projection SL(n)'', and 2025--2026 found
no resolution of this exact rule. The problem retains the source's
generic-data scope.

\subsection{Audit update --- 2026-09-10}\label{audit-update-2026-09-10}

Rechecked
\href{https://pure.tue.nl/ws/files/3846347/391917266748824.pdf}{Baaijens--Draisma},
§4.1, Problem 4.3, and \href{https://arxiv.org/pdf/2501.19310}{Sander's
2025 manuscript}, introduction. The latter still describes the global
selection step as conjectural. Searches for smallest-multiplier
projection results found no general proof. The positive-determinant
projection problem and algorithms reaching stationary points must be
distinguished from the stated \(\det=\pm1\) global rule; impact is
narrowed to specialist.
\end{document}
